\documentclass[11pt,a4paper]{article}
\usepackage[italian]{babel}
\usepackage[margin=2.5cm]{geometry}
\usepackage{parskip}
\usepackage{amsmath, amssymb, amsthm}
\usepackage[most]{tcolorbox}
\usepackage{array}
\usepackage{tikz}
\usetikzlibrary{decorations.pathreplacing}
\usepackage{hyperref}

% --- Riquadri con bordo nero, senza numero ---
% Uso: \begin{definizione} ... \end{definizione}
%      \begin{teorema}[Nome] ... \end{teorema}  (il nome è facoltativo)
\tcbset{nota/.style={colback=white, colframe=black, boxrule=0.5pt,
  sharp corners}}
\NewTColorBox{definizione}{}{nota,
  before upper={\textbf{Definizione.}\ }}
\NewTColorBox{teorema}{o}{nota,
  before upper={\textbf{Teorema\IfValueT{#1}{ (#1)}.}\ }}
\NewTColorBox{proposizione}{o}{nota,
  before upper={\textbf{Proposizione\IfValueT{#1}{ (#1)}.}\ }}
\NewTColorBox{assioma}{o}{nota, boxrule=1pt,
  before upper={\textbf{Assioma\IfValueT{#1}{ (#1)}.}\ }}

% --- Ambienti semplici (senza riquadro) ---
\theoremstyle{definition}
\newtheorem*{esempio}{Esempio}
\newtheorem*{osservazione}{Osservazione}
\newtheorem*{esercizio}{Esercizio}

% V e F nelle tavole di verità
\newcommand{\V}{\textsf{V}}
\newcommand{\F}{\textsf{F}}

\title{Estremi di una funzione e funzioni limitate}
\author{Marco Cerbella}
\date{Lezione 4 -- 5 ottobre 2026}

\begin{document}
\maketitle

\section{Estremo superiore e inferiore di una funzione}

Sia $f : D \to \mathbb{R}$ e $A \subseteq D$.

\begin{definizione}
L'\emph{estremo superiore di $f$ su $A$} è l'estremo superiore dell'insieme $f(A)$:
\[
\sup_{x \in A} f(x) = \sup f(A).
\]
$f$ è \emph{superiormente limitata su $A$} se l'insieme $f(A)$ è superiormente limitato.
\end{definizione}

Se $\sup_{x \in A} f(x)$ è finito e appartiene a $f(A)$, allora è il \emph{valore massimo} di $f$ su $A$. Cioè, se
\begin{enumerate}
    \item per ogni $x \in A$ vale $f(x) \leq M$;
    \item esiste $x_M \in A$ tale che $f(x_M) = M$,
\end{enumerate}
allora $M$ è il valore massimo e $x_M$ è un \emph{punto di massimo}.

\begin{definizione}
L'\emph{estremo inferiore di $f$ su $A$} è l'estremo inferiore dell'insieme $f(A)$:
\[
\inf_{x \in A} f(x) = \inf f(A).
\]
$f$ è \emph{inferiormente limitata su $A$} se $f(A)$ è inferiormente limitato.
\end{definizione}

Se $\inf_{x \in A} f(x)$ è finito e appartiene a $f(A)$, allora è il \emph{valore minimo} di $f$ su $A$. Cioè, se
\begin{enumerate}
    \item per ogni $x \in A$ vale $m \leq f(x)$;
    \item esiste $x_m \in A$ tale che $f(x_m) = m$,
\end{enumerate}
allora $m$ è il valore minimo e $x_m$ è un \emph{punto di minimo}.

\section{Funzioni limitate}

Sia $f : D \to \mathbb{R}$.

\begin{definizione}
\begin{itemize}
    \item $f$ si dice \emph{limitata superiormente} se esiste $M \in \mathbb{R}$ tale che $f(x) \leq M$ per ogni $x \in D$;
    \item $f$ si dice \emph{limitata inferiormente} se esiste $m \in \mathbb{R}$ tale che $f(x) \geq m$ per ogni $x \in D$;
    \item $f$ si dice \emph{limitata} se è limitata sia superiormente sia inferiormente.
\end{itemize}
\end{definizione}

\begin{esempio}
\leavevmode
\begin{enumerate}
    \item $f(x) = x^2$, $D = \mathbb{R}$, $f(\mathbb{R}) = [0, +\infty)$. La funzione è limitata inferiormente ma non superiormente. $m = 0$ è il valore minimo di $f$ e $x_m = 0$ è il punto di minimo.
    \item $f(x) = 3$, $D = \mathbb{R}$, $f(\mathbb{R}) = \{3\}$. La funzione è limitata; $3$ è il valore massimo e il valore minimo, e ogni $x \in \mathbb{R}$ è sia punto di massimo sia punto di minimo.
\end{enumerate}
\end{esempio}

\end{document}
